% FILE: 00_project_identity.tex
% PROJECT: ma — M&A Claim Taxonomy and Board Projection Surface
% THESIS ROLE: ma-claim-taxonomy-and-board-projection-surface

\section{Project Identity}
\label{sec:ma-identity}

\subsection{Purpose}
\label{sec:ma-identity-purpose}

The \texttt{ma/} module defines the complete M\&A-ready process intelligence framework for the
\textit{process-intelligence} research foundry. Its purpose is to establish the formal doctrine
by which process mining evidence is transformed into board-admissible, cryptographically verifiable
claims suitable for use in merger, acquisition, and hostile takeover proceedings.

Conventional M\&A diligence relies on self-reported key performance indicators, static balance
sheets, and PowerPoint presentations whose synergy projections are unverifiable by any buyer
within a practical timeline. The \texttt{ma/} module refutes this regime by defining a
claim taxonomy in which every assertion---from EBITDA optimization to SLA penalty
exposure---must trace through a JSON receipt (BLAKE3-hashed, Ed25519-signed, RFC 8785
JCS-canonical) to an underlying XES or OCEL 2.0 event log and a \texttt{wasm4pm}
conformance query. The governing checkpoint file
\texttt{checkpoint\_\_m\&a-ready\_research\_complete.md} declares this doctrine stack
\textbf{APPROVED} and itemizes all 20 core definition files as Verified.

\subsection{Repository Surface}
\label{sec:ma-identity-surface}

The module occupies \texttt{/Users/sac/process-intelligence/ma/} and contains 40+ specification
documents organized around nine claim taxonomies: Board, Diligence, Synergy, Control, Asset,
Liability, Operational Debt, Integration Risk, and Scalability. Each taxonomy is defined in a
dedicated \texttt{define\_*.md} file cross-referenced by the master consolidation document
\texttt{MASTER\_m\&a-ready\_process\_intelligence\_framework.md} (598 lines, 31,250 bytes).

Five concrete receipt JSON files reside in \texttt{/Users/sac/process-intelligence/receipts/}:

\begin{itemize}
  \item \texttt{rec\_ebitda\_rework\_001.json} --- EBITDA optimization via purchase-order rework
    reduction; fitness 0.982, precision 0.945, EBITDA impact \$1,250,000.
  \item \texttt{rec\_wc\_ar\_002.json} --- Working capital release via accounts-receivable cycle
    compression.
  \item \texttt{rec\_risk\_sla\_003.json} --- SLA penalty exposure via latency trace analysis.
  \item \texttt{rec\_risk\_compliance\_004.json} --- Compliance leakage via LTL rule checking.
  \item \texttt{rec\_residual\_standard\_005.json} --- Process standardization via residual weight
    and entropy analysis.
\end{itemize}

The module also contains the \texttt{M\&A\_BOARD\_PROJECTION\_RENDERER\_SPECIFICATION.md}
(25,258 bytes), which specifies a seven-step pipeline from claims JSON to a cryptographically
sealed PowerPoint deck and Excel workbook, naming specific Rust crates (\texttt{Tera2},
\texttt{pptx-rs}, \texttt{calamine}) and referencing receipt files at concrete absolute paths.

\subsection{Primary Research Role}
\label{sec:ma-identity-role}

Within the dissertation, the \texttt{ma/} module serves as the
\textbf{board-projection surface}: the layer at which the full lineage from 2016 language-model
lessons through the Chatman Equation $A = \mu(O^*)$ through \texttt{wasm4pm} process evidence
is rendered as capital-market-grade financial claims. It demonstrates that process mining is not
merely an operational analytics discipline but a valuation discipline---that every dollar of
merger premium or discount can be grounded in a replayable event log and an independently
verifiable conformance certificate.

The module formalizes the thesis claim that \textit{prediction without coordination is
insufficient}: a board claim that cannot be replayed by an independent buyer within three days
is not a claim---it is narration.

\subsection{Key Artifacts}
\label{sec:ma-identity-artifacts}

\begin{itemize}
  \item \texttt{MASTER\_m\&a-ready\_process\_intelligence\_framework.md} --- Nine-taxonomy
    consolidation with full mathematical formulas.
  \item \texttt{BOARD\_CLAIM\_TAXONOMY.md} --- Six-tier maturity ladder (T1 Existence through
    T6 Recovery) defining evidence requirements and buyer verifiability per tier.
  \item \texttt{REVERSE\_PORTER\_FIVE.md} --- Structural inversion of Porter's Five Forces
    through process evidence moats.
  \item \texttt{hostile-takeover-receipts.md} --- v30.1.1 Hostile Takeover Receipt (HTR)
    taxonomy mandating non-Euclidean, unforgeable process replay in adversarial acquisition
    contexts.
  \item \texttt{SLIDE\_TO\_RECEIPT\_MAP.md} --- Cryptographic mapping schema from slide UUID to
    JSON receipt to event log hash.
  \item \texttt{checkpoint\_\_m\&a-ready\_research\_complete.md} --- APPROVED status checkpoint
    certifying all 20 core definition files mathematically validated.
  \item Five receipt JSON files in \texttt{receipts/} with real BLAKE3 hashes, wasm4pm query
    definitions, fitness/precision numeric results, and Ed25519 validator signatures.
\end{itemize}
